\documentclass[12pt]{article}
\usepackage[margin=0.7in]{geometry}
\usepackage{amsmath}
\usepackage{amssymb}
\usepackage{lmodern}
\usepackage[T1]{fontenc}
\usepackage{microtype}
\pagestyle{empty}

\begin{document}

\begin{center}
{\LARGE Tokens Per Page Expressions}
\end{center}

\vspace{0em}

\begin{align*}
T_{\text{dense}}
&=
\left\lfloor
\frac{P}{H_{\text{kv,loc}} D b}
\right\rfloor
\\[1.2em]
T_{\text{MLA}}
&=
\left\lfloor
\frac{P}{(r_{\text{kv}} + d_{\text{rope}}) b}
\right\rfloor
\end{align*}

\vspace{0.8em}

\noindent\textbf{Where the terms come from}
\begin{itemize}
    \item $P$: physical page size chosen by the runtime. In our experiments, this is fixed at 2 MB.
    \item $H_{\text{kv,loc}}$: local KV-head count, determined by the model architecture and tensor-parallel setup.
    \item $D$: head dimension from the trained model configuration.
    \item $b$: bytes per stored KV-cache element, determined by the cache datatype such as FP16 or BF16.
    \item $r_{\text{kv}}$: MLA KV latent rank from the model configuration.
    \item $d_{\text{rope}}$: MLA RoPE-retained key dimension from the model configuration.
\end{itemize}

\vspace{0.5em}

\noindent\textbf{Main point}

\noindent
The dense and MLA expressions are different, but every term comes from
the model configuration or fixed runtime settings. That is why we can
predict tail waste during inference directly from context length and
known model parameters.

\end{document}
